\section{2014年高考题}
\subsection{2014年大纲卷（理）}\label{gk:2014年大纲卷（理）}


\clearpage


\subsection{2014年大纲卷（文）}\label{gk:2014年大纲卷（文）}


\clearpage


\subsection{2014年全国I卷（理）}\label{gk:2014年全国I卷（理）}


\clearpage


\subsection{2014年全国I卷（文）}\label{gk:2014年全国I卷（文）}


\clearpage


\subsection{2014年全国II卷（理）}\label{gk:2014年全国II卷（理）}


\clearpage


\subsection{2014年全国II卷（文）}\label{gk:2014年全国II卷（文）}


\clearpage


\subsection{2014年安徽卷（理）}\label{gk:2014年安徽卷（理）}


\clearpage


\subsection{2014年安徽卷（文）}\label{gk:2014年安徽卷（文）}


\clearpage


\subsection{2014年北京卷（理）}\label{gk:2014年北京卷（理）}

\begin{description}
    \item[一、]
    \begin{enumerate}
        \item 1
        \item 2
        \item 3
        \item 4
        \item 5
        \item 6
        \item 7
        \item 8
    \end{enumerate}
    \item[二、]
    \begin{enumerate}
      \addtocounter{enumi}{+8}   
        \item 1
        \item 2
        \item 3
        \item 4
        \item 1
        \item 6
    \end{enumerate}
    \item[三、]   
    \begin{enumerate}
      \addtocounter{enumi}{+14}   
        \item 1
        \item 2
        \item 3
        \item 已知函数 $f(x) = x \cos x - \sin x$，$x \in \left[0, \frac{\pi}{2}\right]$。
\begin{enumerate}[label=(\arabic*)]
    \item 求证：$f(x) \leqslant 0$；
    \item 若 $a < \frac{\sin x}{x} < b$ 在 $\left(0, \frac{\pi}{2}\right)$ 上恒成立，求 $a$ 的最大值与 $b$ 的最小值。
\end{enumerate}
        \item 已知椭圆 $C: x^{2} + 2y^{2} = 4$。
\begin{enumerate}[label=(\arabic*)]
    \item 求椭圆 $C$ 的离心率；
    \item 设 $O$ 为坐标原点，若点 $A$ 在椭圆 $C$ 上，点 $B$ 在直线 $y = 2$ 上，且 $OA \perp OB$，试判断直线 $AB$ 与圆 $x^{2} + y^{2} = 2$ 的位置关系，并证明你的结论。
\end{enumerate}
        \item 对于数对序列 $P: (a_{1}, b_{1}), (a_{2}, b_{2}), \cdots, (a_{n}, b_{n})$，
        记 $T_{1}(P) = a_{1} + b_{1}$，
        $T_{k}(P) = b_{k}$ + $\max \left\{ T_{k-1}(P), \right.$ $\left. a_{1} + a_{2} + \cdots + a_{k} \right\} \quad (2 \leqslant k \leqslant n)$，其中 $\max \left\{ T_{k-1}(P), a_{1} + a_{2} + \cdots + a_{k} \right\}$ 表示 $T_{k-1}(P)$ 和 $a_{1} + a_{2} + \cdots + a_{k}$ 两个数中最大的数。

\begin{enumerate}[label=(\arabic*)]
    \item 对于数对序列 $P: (2,5), (4,1)$，求 $T_{1}(P)$, $T_{2}(P)$ 的值；
    \item 记 $m$ 为 $a, b, c, d$ 四个数中最小值，对于由两个数对 $(a, b)$, $(c, d)$ 组成的数对序列 $P: (a, b), (c, d)$ 和 $P': (c, d), (a, b)$，试分别对 $m = a$ 和 $m = d$ 时两种情况比较 $T_{2}(P)$ 和 $T_{2}(P')$ 的大小；
    \item 在由 5 个数对 $(11,8), (5,2), (16,11), (11,11), (4,6)$ 组成的所有数对序列中，写出一个数对序列 $P$ 使 $T_{5}(P)$ 最小，并写出 $T_{5}(P)$ 的值。（只需写出结论）
\end{enumerate}
    \end{enumerate}
\end{description}
\clearpage


\subsection{2014年北京卷（文）}\label{gk:2014年北京卷（文）}

\begin{description}
    \item[一、]
    \begin{enumerate}
        \item 1
        \item 2
        \item 3
        \item 4
        \item 5
        \item 6
        \item 7
        \item 8
    \end{enumerate}
    \item[二、]
    \begin{enumerate}
      \addtocounter{enumi}{+8}   
        \item 1
        \item 2
        \item 3
        \item 4
        \item 1
        \item 6
    \end{enumerate}
    \item[三、]   
    \begin{enumerate}
      \addtocounter{enumi}{+14}   
        \item 1
        \item 2
        \item 3
        \item 
        \item 已知椭圆 $C: x^{2} + 2y^{2} = 4$。
\begin{enumerate}[label=(\arabic*)]
    \item 求椭圆 $C$ 的离心率；
    \item 设 $O$ 为原点，若点 $A$ 在直线 $y = 2$ 上，点 $B$ 在椭圆 $C$ 上，且 $OA \perp OB$，求线段 $AB$ 长度的最小值。
\end{enumerate}
        \item 已知函数 $f(x) = 2x^{3} - 3x$。
\begin{enumerate}[label=(\arabic*)]
    \item 求 $f(x)$ 在区间 $[-2,1]$ 上的最大值；
    \item 若过点 $P(1,t)$ 存在 3 条直线与曲线 $y = f(x)$ 相切，求 $t$ 的取值范围；
    \item 问过点 $A(-1,2)$, $B(2,10)$, $C(0,2)$ 分别存在几条直线与曲线 $y = f(x)$ 相切？（只需写出结论）
\end{enumerate}
    \end{enumerate}
\end{description}
\clearpage


\subsection{2014年重庆卷（理）}\label{gk:2014年重庆卷（理）}


\clearpage


\subsection{2014年重庆卷（文）}\label{gk:2014年重庆卷（文）}


\clearpage


\subsection{2014年福建卷（理）}\label{gk:2014年福建卷（理）}


\clearpage


\subsection{2014年福建卷（文）}\label{gk:2014年福建卷（文）}


\clearpage


\subsection{2014年广东卷（理）}\label{gk:2014年广东卷（理）}


\clearpage


\subsection{2014年广东卷（文）}\label{gk:2014年广东卷（文）}


\clearpage


\subsection{2014年湖北卷（理）}\label{gk:2014年湖北卷（理）}


\clearpage


\subsection{2014年湖北卷（文）}\label{gk:2014年湖北卷（文）}


\clearpage


\subsection{2014年湖南卷（理）}\label{gk:2014年湖南卷（理）}


\clearpage


\subsection{2014年湖南卷（文）}\label{gk:2014年湖南卷（文）}


\clearpage


\subsection{2014年江苏卷}\label{gk:2014年江苏卷}


\clearpage


\subsection{2014年江西卷（理）}\label{gk:2014年江西卷（理）}


\clearpage


\subsection{2014年江西卷（文）}\label{gk:2014年江西卷（文）}


\clearpage


\subsection{2014年辽宁卷（理）}\label{gk:2014年辽宁卷（理）}


\clearpage


\subsection{2014年辽宁卷（文）}\label{gk:2014年辽宁卷（文）}


\clearpage


\subsection{2014年上海卷（理）}\label{gk:2014年上海卷（理）}


\clearpage


\subsection{2014年上海卷（文）}\label{gk:2014年上海卷（文）}


\clearpage


\subsection{2014年四川卷（理）}\label{gk:2014年四川卷（理）}


\clearpage


\subsection{2014年四川卷（文）}\label{gk:2014年四川卷（文）}


\clearpage


\subsection{2014年山东卷（理）}\label{gk:2014年山东卷（理）}


\clearpage


\subsection{2014年山东卷（文）}\label{gk:2014年山东卷（文）}


\clearpage


\subsection{2014年陕西卷（理）}\label{gk:2014年陕西卷（理）}


\clearpage


\subsection{2014年陕西卷（文）}\label{gk:2014年陕西卷（文）}


\clearpage


\subsection{2014年天津卷（理）}\label{gk:2014年天津卷（理）}


\clearpage


\subsection{2014年天津卷（文）}\label{gk:2014年天津卷（文）}


\clearpage


\subsection{2014年浙江卷（理）}\label{gk:2014年浙江卷（理）}


\clearpage


\subsection{2014年浙江卷（文）}\label{gk:2014年浙江卷（文）}


\clearpage

